\documentclass[11pt]{article}

% ---------------------------------------------------------------------------
%  L5 — Convex king animals.
%  Category L (entirely machine-written); see paper/README.md and
%  docs/publication-split.md.
%
%  N3 COLLISION.  Gouyou-Beauchamps & Leroux, FPSAC 2004, §2.3, have the block
%  decomposition of Lemma 1 and the mirror equality of Proposition 9, for
%  convex polyominoes on the honeycomb lattice.  docs/publication-split.md
%  requires that citation in THREE places: at Lemma 1, at Proposition 9, and
%  in the related-work section.  Do not let a revision drop one of them.
%
%  Source material, all in this repository:
%    results/hv-growth-sandwich.md      Props 6-11, Lemmas 1-4, Conjecture 8
%    results/convex-polyplets.md        the exclusions, mu, the PSLQ boxes
%    docs/proofs/convex-mirage.md       the original mirage finding
%    results/middle-kingdom-phase3.md   the grid, Corollary 4, Proposition 5
%    results/middle-kingdom.md          the campaign index
%    results/mk-dir4-perimeter.md       the dir4 perimeter series
% ---------------------------------------------------------------------------

% ---------------------------------------------------------------------------
%  paper/shared/preamble.tex — packages and macros common to every manuscript
%  in this directory.  Factored out of paper/polyplets-report.tex and
%  paper/technical-report.tex, which between them had two copies of the same
%  twenty lines.
%
%  technical-report.tex does NOT \input this file and is not to be edited to do
%  so: it is jasonp's prose and is read-only to the machine (paper/README.md,
%  "Who may edit what").  The duplication there is deliberate.
%
%  Assumes \documentclass[11pt]{article} has already been issued.
% ---------------------------------------------------------------------------
\usepackage[margin=1in]{geometry}
\usepackage{amsmath,amssymb}
\usepackage{amsthm}
\usepackage{booktabs}
\usepackage{graphicx}
\usepackage{numprint}
\usepackage{microtype}
\usepackage[table]{xcolor}
\usepackage[hidelinks]{hyperref}
\usepackage{flafter}   % a float never appears before the text that introduces it
\usepackage{float}     % [H] pins tables/figures exactly where written
\usepackage[font=small,labelfont=bf,labelsep=period]{caption}
\setlength{\intextsep}{16pt plus 3pt minus 3pt}

\npthousandsep{,}
\npfourdigitsep
\newcommand{\num}[1]{\numprint{#1}}
\newcommand{\g}{\cellcolor{gray!15}}

% Theorem environments.  One counter shared across kinds, so that a reference
% like "Theorem 4" is unambiguous in a paper that also has lemmas and
% propositions.
\theoremstyle{plain}
\newtheorem{theorem}{Theorem}
\newtheorem{proposition}[theorem]{Proposition}
\newtheorem{lemma}[theorem]{Lemma}
\newtheorem{corollary}[theorem]{Corollary}
\theoremstyle{definition}
\newtheorem{definition}[theorem]{Definition}
\newtheorem{example}[theorem]{Example}
\theoremstyle{remark}
\newtheorem{remark}[theorem]{Remark}
\newtheorem{conjecture}[theorem]{Conjecture}
\newtheorem{openproblem}[theorem]{Open Problem}

% Repo-path citations.  Every one of these must resolve in the tree that ships
% with the paper; `make gate-citations` checks the markdown, and
% paper/verify_claims.py checks the manuscripts.
\newcommand{\repo}[1]{\texttt{#1}}

% OEIS links.
\newcommand{\oeis}[1]{\href{https://oeis.org/#1}{#1}}

% Growth constants, so a rename is one edit.
\newcommand{\lam}{\lambda}
\newcommand{\muH}{\mu_H}

% ---------------------------------------------------------------------------
%  paper/shared/disclosure.tex — the two authorship disclosure blocks.
%
%  docs/publication-split.md §1 fixes the rule these implement:
%
%    "Under no circumstances will there be any ambiguity about whether a paper
%     is wholly my work (none), merely my prose (some) or entirely LLM-authored
%     (some)."   — jasonp, 2026-08-07
%
%  There are exactly two blocks and they live in exactly one file, so that a
%  paper cannot quietly soften its own disclosure.  The fixed sentences take no
%  arguments.  The ONE thing a paper supplies is the per-result verification
%  ledger, because that is the part that differs honestly between papers — and
%  §1 requires it to be per result, not in aggregate.
%
%  Do not add a \Pdisclosure or \Ldisclosure variant with an optional softening
%  argument.  If a paper does not fit one of these two blocks, the paper is
%  wrong, not the block.
% ---------------------------------------------------------------------------

\newsavebox{\disclosurebox}

\newenvironment{disclosureframe}
  {\begin{center}\begin{lrbox}{\disclosurebox}\begin{minipage}{0.93\textwidth}\small}
  {\end{minipage}\end{lrbox}\fbox{\usebox{\disclosurebox}}\end{center}}

% --- P: jasonp's prose -------------------------------------------------------
% Byline: Jason Parker, alone.  Argument: the per-result ledger of what the
% machine did.
\newcommand{\Pdisclosure}[1]{%
\begin{disclosureframe}
\textbf{Authorship.}\enspace Every sentence of this paper was written by its
named author. He has read every proof it states and believes each one, and he
is responsible for the correctness of what it claims.

\smallskip
\textbf{Machine contribution.}\enspace #1
\end{disclosureframe}}

% --- L: entirely machine-written --------------------------------------------
% Argument: the per-result verification ledger — for each result, what warrants
% it (Lean kernel, exact certificate, or labelled measurement) and how much of
% it a human has checked.
\newcommand{\Ldisclosure}[1]{%
\begin{disclosureframe}
\textbf{Authorship.}\enspace This document was written end to end by a large
language model. That includes its mathematics, its prose, its tables and its
\LaTeX{}. No sentence of it was written by a human. The mathematics it reports
was likewise derived by machine.

\smallskip
\textbf{Human role.}\enspace The person who directed this work chose what to
compute, chose what to publish, and is responsible for the decision to release
this document. Except where the ledger below says otherwise, that person has
not personally checked the arguments here, and this disclosure is not to be
read as an endorsement of them.

\smallskip
\textbf{What warrants each result.}\enspace #1
\end{disclosureframe}}

% --- Draft banner ------------------------------------------------------------
% For an L-paper whose verification ledger is not yet settled.  Loud on purpose:
% an L-paper without a truthful ledger must not be mistaken for a finished one.
%
% \firstdraftbanner is the standard wording, which five papers previously
% carried character-for-character in their own preambles.  It lives here for the
% same reason the disclosure blocks do: identical text in seven places drifts,
% and a paper must not be able to soften it by editing its own copy.  The
% optional argument is for a gate specific to one paper, appended verbatim; a
% paper needing different wording (L6, compute-gated) calls \draftbanner direct.
\newcommand{\firstdraftbanner}[1][]{%
\draftbanner{This is a first draft. The verification ledger below records what
warrants each result \emph{mechanically}; the column that records what a human
has read is empty, deliberately, because nobody has read it yet. Do not
circulate until that column says something true.#1}}

\newcommand{\draftbanner}[1]{%
\begin{center}
\fcolorbox{red!60!black}{yellow!15}{%
  \begin{minipage}{0.93\textwidth}\small
  \textbf{DRAFT — NOT FOR CIRCULATION.}\enspace #1
  \end{minipage}}
\end{center}}


\newcommand{\Z}{\mathbb{Z}}
\newcommand{\Q}{\mathbb{Q}}
\newcommand{\R}{\mathbb{R}}
\newcommand{\HV}{\mathrm{HV}}

\title{\textbf{Convex king animals}\\[4pt]
\large Area is wild, perimeter is tame, and one growth constant
serves every class in between}
\author{[byline: jasonp's call --- \repo{docs/publication-split.md} \S1]}
\date{August 2026}

\begin{document}
\maketitle

\draftbanner{This is a first draft. The verification ledger below records what
warrants each result \emph{mechanically}; the column that records what a human
has read is empty, deliberately, because nobody has read it yet. Do not
circulate until that column says something true. One further gate is specific to
this paper: the attribution to Gouyou-Beauchamps and Leroux
(\S\ref{sec:related}) must survive every revision, in all three of the places it
appears.}

\Ldisclosure{%
\emph{Proposition~\ref{prop:squeeze} (the squeeze) and Lemmas~\ref{lem:blocks},
\ref{lem:stacks}, \ref{lem:supermul}} --- proved below, with no analytic input:
an injection, Fekete's lemma, and an elementary partition bound.
Lemma~\ref{lem:supermul}'s column join and its inverse are formalized in Lean~4
against mathlib (\texttt{Stair.join\_valid}, \texttt{Stair.cut\_join},
\texttt{Stair.join\_injOn}), with axiom footprint \texttt{[propext,
Quot.sound]} --- not even \texttt{Classical.choice} --- and the counting layer
and Fekete step are there too (\texttt{Stair.M\_supermul},
\texttt{Stair.M\_tendsto}, \texttt{Stair.M\_le\_mu\_pow}), all guarded. So
Lemma~\ref{lem:supermul} is a Lean theorem.
\emph{Proposition~\ref{prop:squeeze} itself is NOT formalized}: the geometric
HV-convex layer (Lemma~\ref{lem:blocks}) and the squeeze are paper proofs.
\emph{The exclusions of \S\ref{sec:exclusions}} --- rigorous in the direction
claimed: full column rank modulo a prime is a proof of full rank over $\Q$.
Four positive and negative control arms plus a null control are pinned by a
gate.
\emph{$\mu$ to 199 digits, $\theta = 0$, the amplitudes} --- labelled
measurement, with the trusted digit count itself measured rather than asserted.
\emph{The PSLQ verdicts} --- labelled measurement, stated as falsifiable box
exclusions and never as ``$\mu$ is transcendental''.
\emph{Lemma~\ref{lem:T} and Proposition~\ref{prop:split}} --- proved below.
\emph{Conjecture~\ref{conj:sharp}} --- \textbf{a conjecture}, and the paper is
explicit that the repository has no proof of a sharp asymptotic for \emph{any}
series in this family.
\emph{Propositions~\ref{prop:mirror}, \ref{prop:phi}, \ref{prop:ratio}} ---
proved; Proposition~\ref{prop:ratio} is conditional on the amplitudes existing,
which is strictly weaker than Conjecture~\ref{conj:sharp}.
\emph{Novelty} --- N1 and N2 clean; \textbf{N3 is a real collision}, cited in
three places.
\emph{Human verification:} none, as of this draft.}

\begin{abstract}
\noindent
A \emph{convex king animal} is a king-connected set of cells of $\Z^2$ in which
every row and every column is a single contiguous run. We report what is now
known about counting them.

By \textbf{area} the class is intractable, and we make that quantitative. On
$700$ terms the generating function satisfies no P-recurrence of order $\le 24$
with polynomial coefficients of degree $\le 24$, and no algebraic equation of
degree $\le 20$ in $F$ with coefficients of degree $\le 20$ in $t$; the same
holds for the classical convex-polyomino control regenerated to the same length.
The growth constant is
$\mu = 3.128943269730886\dots$, pinned to $199$ digits, with the correction to
pure exponential growth measured to be geometric --- $\theta = 0$, no power-law
factor and nothing resembling a stretched exponential. Integer-relation searches
find no polynomial with $\mu$ as a root of degree $\le 20$ and height $\le 10^8$,
$\le 30$ and $\le 10^6$, or $\le 12$ and $\le 10^{15}$.

By \textbf{semiperimeter} the same class is algebraic, of degree $2$; imposing a
directedness condition strong enough to cut the population to about $1\%$ raises
the degree to $4$ and does not leave the algebraic class. Convexity is a
perimeter lever, not an area lever --- for king animals exactly as for
polyominoes.

The structural result is a squeeze. Let $\mathcal C$ be \emph{any} class with
$\text{staircase} \subseteq \mathcal C \subseteq \text{HV-convex}$. Then
$\lim \mathcal C(n)^{1/n}$ exists and equals $\mu$. Reading an HV-convex animal
left to right, it fattens, shears, then thins; the shearing middle is a
staircase animal and the two ends are stacks, which are too few to move an
exponential. The proof uses no analytic input and, because the staircase join is
exactly injective, makes the limit a supremum --- so every enumerated term is a
rigorous floor under $\mu$.

Finally we locate the subdominant structure. The four-cone-directed subclass
splits exactly into an ascending half growing at $\mu$ and a descending half
growing at a second constant $\nu = 2.5145796\dots$, both proved; and the
amplitude ratio of the two classes, previously known only as $54$ measured
digits, has a closed expression as a ratio of two explicit vectors against the
bounded solution of a three-term recurrence. The same recurrence computes $\mu$
by shooting --- $987$ digits in $1.6$ seconds, with no series and no
extrapolation.
\end{abstract}

\tableofcontents

\section{Introduction}

Convexity is the classical tractability lever for polyominoes. Convex polygons
by perimeter have an algebraic generating
function~\cite{delestViennot1984}; convex polyominoes by area grow at
$2.30914\dots$~\cite{bender1974}; symmetry classes of convex polyominoes have
been enumerated on the square lattice~\cite{lrr1998} and on the honeycomb
lattice~\cite{gblRoux2004}. It is natural to expect the lever to work on the
king lattice too.

It does, but only for perimeter. That split --- \emph{area wild, perimeter tame}
--- is the organising fact of this paper, and it holds for both lattices, which means the classical solvability results were always
perimeter results and the area side was never tractable.

\begin{definition}
A \emph{convex king animal}, equivalently an \emph{HV-convex king animal} or
\emph{convex polyplet}, is a finite king-connected set of cells of $\Z^2$ in
which every row and every column is a single contiguous run, counted up to
translation. Write $A(n)$ for the number with $n$ cells:
\[
  1,\; 4,\; 16,\; 61,\; 221,\; 766,\; 2566,\; 8390,\; 26982,\; 85834,\;
  271174,\; 853111,\; \dots
\]
A \emph{staircase} king animal is one in which, scanning columns left to right,
both the bottom and the top of the column interval are nondecreasing. Write
$M(n)$ for their number: $1, 3, 9, 28, 87, 272, \dots$ (OEIS \oeis{A225114}, the skew shapes with no empty row or column).
\end{definition}

\section{The class and its transfer structure}
\label{sec:structure}

An HV-convex king animal is a sequence of column intervals $[b(j), t(j)]$,
$j = 1,\dots,k$, with $b$ valley-unimodal and $t$ peak-unimodal, counted up to
translation. Write
\[
  d(j) = b(j{+}1) - b(j), \qquad h(j) = t(j) - b(j) + 1, \qquad
  s(j) = h(j) - h(j{+}1).
\]

Give column $j$ the \emph{phase} $(p_b, p_t)$, where $p_b = 1$ if some earlier
step had $d > 0$ and $p_t = 1$ if some earlier step had $t$ decreasing.
Unimodality makes both bits monotone in $j$: each flips at most once and never
flips back. So the transfer operator is block-triangular in the four phases, and
the diagonal blocks are:

\begin{table}[H]
\centering\small
\begin{tabular}{c l l l}
\toprule
phase & meaning & step constraint & kernel $K(h,h')$ \\
\midrule
$(0,0)$ & $b$ falls, $t$ rises & $s \le d \le 0$ & $h'-h+1$, \; $h' \ge h$ \\
$(1,0)$ & both rise & $d \ge \max(0,s)$ & $\min(h,h')+1$ \\
$(0,1)$ & both fall & $d \le \min(0,s)$ & $\min(h,h')+1$ \\
$(1,1)$ & $b$ rises, $t$ falls & $0 \le d \le s$ & $h-h'+1$, \; $h' \le h$ \\
\bottomrule
\end{tabular}
\caption{The four diagonal blocks. Blocks $(1,0)$ and $(0,1)$ are the staircase
kernel and its mirror; blocks $(0,0)$ and $(1,1)$ are stacks. The kernel
$\min(h,h')+1$ is the whole of the king lattice's difference from the square
one, where it is $\min(h,h')$.}
\label{tab:blocks}
\end{table}

The unbounded state variable is the column height, and the coupling in it is
the single kernel $\min(h,h')+1$, appearing once in each middle block. Almost
everything below is an argument about that one kernel.

\section{The squeeze}
\label{sec:squeeze}

\begin{proposition}[the squeeze]
\label{prop:squeeze}
Let $\mathcal C$ be any class of king animals with
\[
  \text{staircase} \;\subseteq\; \mathcal C \;\subseteq\; \text{HV-convex}.
\]
Then $\lim_n \mathcal C(n)^{1/n}$ exists and equals
$\lim_n A(n)^{1/n} = \mu = 3.128943269730886\dots$
\end{proposition}

Read informally: $A(n)$ counts blobs with no gap in any row or column; $M(n)$
counts the staircase ones, a thin slice. The claim is that they grow at the same
exponential rate, as does everything between. Every intermediate class is
trapped between $M$ and $A$ term by term, so nothing has to be proved about the individual classes.

\begin{lemma}[three-block factorisation]
\label{lem:blocks}
Write $P(n)$ for the number of monotone-height blocks of area $n$ (phase
$(0,0)$, or equivalently by reversal phase $(1,1)$) and $M(n)$ for the staircase
count, with $P(0) = M(0) = 1$. Then
\begin{equation}
  A(n) \;\le\; 2\,(n{+}1)^2 \sum_{i+j+l=n} P(i)\,M(j)\,P(l) .
  \label{eq:factor}
\end{equation}
\end{lemma}

\begin{proof}
By the monotonicity of the phase bits, the columns split into three consecutive
(possibly empty) runs $C_1C_2C_3$ carrying phases $(0,0)$; one of $(1,0)$ or
$(0,1)$; and $(1,1)$. Every step internal to a run obeys that run's constraint
in Table~\ref{tab:blocks}, so taken standalone $C_1$ and $C_3$ are
monotone-height blocks and $C_2$ is a staircase animal or the vertical mirror of
one. The animal is recovered from $(C_1,C_2,C_3)$, the two junction offsets and
one bit naming $C_2$'s phase; each junction offset lies in $[-h', h]$ with
$h + h' \le n$, hence takes at most $n+1$ values.
\end{proof}

\begin{remark}[attribution --- 1 of 3]
\label{rem:attr1}
This decomposition is not new in kind. Gouyou-Beauchamps and
Leroux~\cite{gblRoux2004}, \S2.3 ``Growth phases of convex polyominoes'',
decompose a convex polyomino into blocks $H_{ij}$ by the growth phases of its
upper and lower profiles, with the column state an ordered pair, the transitions
one-way, the extreme blocks $H_{00}$ and $H_{22}$ identified as \emph{stack}
polyominoes and the middle blocks as \emph{staircase} polyominoes with
$H_{02} = Pa = H_{20}$. That is Lemma~\ref{lem:blocks} here, together with
Lemma~\ref{lem:stacks}'s identification and Proposition~\ref{prop:mirror}, for
convex polyominoes on the honeycomb lattice. Their square-lattice
companion~\cite{lrr1998} reaches the same objects by a different route, so 2004
is the source to cite. What is not theirs: the king lattice --- their middle
kernel is $\min(h,h')$ and ours is $\min(h,h')+1$ --- a class with no exact
solution, and everything Propositions~\ref{prop:squeeze}, \ref{prop:split},
\ref{prop:phi} and~\ref{prop:ratio} do with the blocks.
\end{remark}

\begin{lemma}[the outer blocks are stacks, hence sub-exponential]
\label{lem:stacks}
$P(n)$ is the number of weakly unimodal compositions of $n$ (OEIS
\oeis{A001523}), and $P(n) \le E(n) := (n{+}1)^{4\sqrt n + 6}$, so
$P(n)^{1/n} \to 1$.
\end{lemma}

The lemma is two claims bolted together --- \emph{what} the outer blocks are (a
bijection) and \emph{how few} there are (a counting bound). Only the second is
used downstream.

\begin{proof}[Part 1: a stack is a weakly unimodal composition]
In a phase-$(1,1)$ block the bottoms rise and the tops fall, so the column
intervals are \emph{nested}: $[b(j{+}1),t(j{+}1)] \subseteq [b(j),t(j)]$.
Nesting means that the set of columns meeting row $y$ is never broken --- it is
a prefix $1..r(y)$. Transpose the picture and record the row widths $r(y)$ from
the bottom up. Since $\{y : r(y) \ge k\} = [b(k),t(k)]$ is a decreasing nested
family of intervals, $r$ rises then falls, and $\sum_y r(y) = n$. The stack is
rebuilt from $r$ alone, so the map is a bijection.

Worked instance: columns $[0,4],[1,3],[2,2]$ (heights $5,3,1$, area $9$) have
row widths $r = (1,2,3,2,1)$, unimodal and summing to $9$.
\end{proof}

\begin{proof}[Part 2: there are few of them]
Cut a weakly unimodal composition at its peak; what falls out on either side is
a partition, so $P(n) \le (n{+}1)^2 p(n)^2$, the $(n{+}1)^2$ paying for the
peak's position and value.

For $p(n)$ a deliberately crude bound suffices, since the squeeze uses only
$P(n)^{1/n}\to1$. Put $s = \lceil\sqrt n\rceil$ and split a partition of $n$ at
$s$. The parts $\le s$ are determined by their $s$ multiplicities, each an
integer in $[0,n]$, contributing at most $(n{+}1)^s$ choices. The parts $> s$
number at most $L = \lfloor n/(s{+}1)\rfloor$, and listing them in nonincreasing
order gives at most $(L{+}1)n^L \le (n{+}1)^{L+1}$ choices. Hence
$p(n) \le (n{+}1)^{s+L+1} \le (n{+}1)^{2\sqrt n + 2}$, since $s \le \sqrt n + 1$
and $L \le \sqrt n$. So $P(n) \le E(n)$ and
$\log P(n)/n \le (4\sqrt n + 6)\log(n{+}1)/n \to 0$.
\end{proof}

Only the exponent matters: $\sqrt n \log n$ is $o(n)$, so however many stacks
there are, they carry no exponential. Hardy and
Ramanujan~\cite{hardyRamanujan1918} would give the sharper $p(n) = e^{O(\sqrt
n)}$; nothing downstream distinguishes the two, and the elementary bound keeps
the squeeze free of analytic input, which is worth more here than sharpness.

\begin{lemma}[the staircase count is supermultiplicative]
\label{lem:supermul}
$M(i)M(j) \le M(i+j)$ for all $i,j \ge 0$. Hence $\lim_n M(n)^{1/n}$ exists and
equals $\sup_n M(n)^{1/n}$.
\end{lemma}

\begin{proof}
Given staircase animals $X$ of area $i$ with last column height $h$, and $Y$ of
area $j$ with first column height $h'$, slide $Y$ up by $d = \max(0, h-h')$ and
butt the two column sequences together.

\emph{Why that $d$ and no other.} Staircase means both boundaries
nondecreasing. At the seam, $b$ not dropping forces $d \ge 0$; $t$ not dropping
forces $d \ge h - h'$. So $\max(0, h-h')$ is the smallest legal slide --- it is a
\emph{function} of $X$ and $Y$, not a choice, which is what stops the join from
smuggling in information.

\emph{The join is in the class.} $d \ge 0$ keeps bottoms nondecreasing;
$d - s = \max(0,h-h') - (h-h') \ge 0$ keeps tops nondecreasing; and $d \le h$,
since $d = h-h' \le h-1$ when $h \ge h'$ and $d = 0$ otherwise, so the junction
columns are king-adjacent. Area is $i+j$.

\emph{The join is injective at fixed $(i,j)$.} Every column is nonempty, so the
cumulative areas $h(1)+\dots+h(r)$ are \emph{strictly} increasing, and a
strictly increasing sequence hits $i$ at most once. Cut at that unique prefix:
those are $X$'s columns and the rest are $Y$'s, each read back up to
translation, which is how both were counted. No split index has to be supplied
alongside the join, so distinct pairs have distinct images.

Fekete's lemma in supermultiplicative form finishes.
\end{proof}

\begin{corollary}[every banked term is a floor]
\label{cor:floor}
$\mu \ge M(n)^{1/n}$ for every $n$, with no extrapolation. At $n = 700$,
\[
  \mu \;\ge\; M(700)^{1/700} \;=\; 3.12340450886853853211\dots
\]
\end{corollary}

The digits past $3.12$ in the extrapolated value remain extrapolation; the
inequality above is exact arithmetic on an enumerated term.

\begin{proof}[Proof of Proposition~\ref{prop:squeeze}]
Write $\mu = \lim M(n)^{1/n} = \sup_n M(n)^{1/n} \ge 1$, which exists by
Lemma~\ref{lem:supermul}. Fix $\varepsilon > 0$ and take $C_\varepsilon$ with
$M(j) \le C_\varepsilon(\mu+\varepsilon)^j$. In \eqref{eq:factor}, bound
$P(i), P(l) \le E(n)$ by Lemma~\ref{lem:stacks}, valid for every $i, l \le n$
since $E$ is nondecreasing:
\[
  A(n) \;\le\; 2(n{+}1)^4 E(n)^2 C_\varepsilon (\mu+\varepsilon)^n ,
\]
and $E(n)^{1/n} \to 1$, so $\limsup A(n)^{1/n} \le \mu + \varepsilon$ for every
$\varepsilon$. In the other direction, staircase $\subseteq \mathcal C
\subseteq$ HV-convex gives $M(n) \le \mathcal C(n) \le A(n)$ termwise, so
$\liminf \mathcal C(n)^{1/n} \ge \mu$ and
$\limsup \mathcal C(n)^{1/n} \le \limsup A(n)^{1/n} \le \mu$.
\end{proof}

\subsection{What the squeeze settles}

The four-cone-directed subclass shares $\mu$ with the unrestricted one to $49$
digits, which had looked like a bijection waiting to be found. It is not: it is
Proposition~\ref{prop:squeeze}, and the four-cone condition is not special.
\emph{Every} condition that leaves the staircase animals alone leaves $\mu$
alone, because the staircase animals already carry all of the exponential
entropy. Table~\ref{tab:blocks} says the same thing at the operator level: the
four-cone floor rewrites the kernels of blocks $(0,0)$ and $(0,1)$ and leaves
blocks $(1,0)$ and $(1,1)$ untouched entry for entry --- and block $(1,0)$ is the
staircase kernel.

\begin{table}[H]
\centering\small
\begin{tabular}{l l r l}
\toprule
class & first terms & trusted digits of $\mu$ & amplitude $C$ \\
\midrule
staircase (\oeis{A225114}) & $1,3,9,28,87,272$ & 204 & $0.28932146397165132308$ \\
HV-convex, $b$ nondecreasing & $1,3,10,33,107,342$ & 202 & $0.37545302027992473174$ \\
(dir4, HV-convex) & $1,4,15,53,177,567$ & --- & $0.45030318571234118235$ \\
HV-convex & $1,4,16,61,221,766$ & 199 & $0.97445221313500464915$ \\
\bottomrule
\end{tabular}
\caption{Four classes, one constant. The staircase value at its $204$ trusted
digits reproduces \emph{every one} of the $199$ the HV-convex series pins and
continues past them. Three independent series, three different classes.}
\label{tab:fourrungs}
\end{table}

\section{Area is wild}
\label{sec:exclusions}

\begin{theorem}
\label{thm:excl}
On $700$ terms, the by-area generating function of HV-convex king animals
satisfies no P-recurrence
$\sum_{i \le J} p_i(n)a(n{+}i) = 0$ with $J \le 24$ and $\deg p_i \le 24$, and no
algebraic equation $\sum_{j \le K} q_j(t)F(t)^j = 0$ with $K \le 24$ and
$\deg q_j \le 24$. The same holds for HV-convex polyominoes by area
(\oeis{A067675}), regenerated to the same length.
\end{theorem}

Two things make this one test rather than a $25\times25$ sweep, and they are what makes the negative rigorous.

\paragraph{The boxes nest.} A solution with $J' \le J$ and $D' \le D$ is a
solution of the $(J,D)$ ansatz with the unused $p_i$ set to zero. Excluding the
maximal box excludes every box inside it.

\paragraph{Full column rank mod $p$ is a proof over $\Q$.} The matrix entries
are integers. Full column rank modulo $p$ means some maximal minor is nonzero
mod $p$, hence nonzero over $\Q$, hence the only rational solution is trivial.
Rank can drop mod $p$, never rise, so \emph{this} direction is rigorous. A rank
\emph{defect} would only be a candidate --- which is the branch that gets the
holdout and the second prime.

\begin{table}[H]
\centering\small
\begin{tabular}{l l l r r r l}
\toprule
series & mode & box & unknowns & rows & rank & verdict \\
\midrule
king    & prec & $(20,20)$ & 441 & 680 & 441 & EXCLUDED \\
king    & prec & $(24,24)$ & 625 & 676 & 625 & EXCLUDED \\
king    & prec & $(20,20)$, skip 50 rows & 441 & 630 & 441 & EXCLUDED \\
king    & prec & $(20,20)$, second prime & 441 & 680 & 441 & EXCLUDED \\
king    & alg  & $(20,20)$ & 441 & 701 & 441 & EXCLUDED \\
king    & alg  & $(24,24)$ & 625 & 701 & 625 & EXCLUDED \\
control & prec & $(20,20)$ & 441 & 680 & 441 & EXCLUDED \\
control & prec & $(24,24)$ & 625 & 676 & 625 & EXCLUDED \\
control & prec & $(20,20)$, skip 50 rows & 441 & 630 & 441 & EXCLUDED \\
control & prec & $(20,20)$, second prime & 441 & 680 & 441 & EXCLUDED \\
control & alg  & $(20,20)$ & 441 & 701 & 441 & EXCLUDED \\
control & alg  & $(24,24)$ & 625 & 701 & 625 & EXCLUDED \\
\bottomrule
\end{tabular}
\caption{$p = 2^{61}-1$ unless noted. $(25,25)$ is the first box that $700$
terms cannot decide --- $676$ unknowns against $675$ rows --- so that is where
the term count runs out, not the method. The row skip tests ``holds eventually''
rather than ``holds from $n=1$''; same verdict.}
\label{tab:excl}
\end{table}

Algebraicity was tested separately even though algebraic implies D-finite, since
the implication only rules out algebraic functions whose induced ODE fits inside
the box already excluded, and a degree-$24$ algebraic function can induce a much
larger one.

\subsection{The guesser is powered}

A negative from a guesser means nothing without a positive control. Four arms
run as a gate, each required to come out as stated:

\begin{table}[H]
\centering\small
\begin{tabular}{l l l l l}
\toprule
arm & series & box & must be & is \\
\midrule
prec $+$ & HV-convex by semiperimeter & $(5,2)$ & CANDIDATE & nullity 1; predicts all 172 held-out rows \\
prec $-$ & same series & $(2,1)$ & EXCLUDED & rank 6 of 6 \\
alg $+$ & \oeis{A005436}, algebraic~\cite{delestViennot1984} & $(2,8)$ & CANDIDATE & nullity 1; holds on all 69 held-out rows \\
alg $-$ & same series & $(2,4)$ & EXCLUDED & rank 15 of 15 \\
\bottomrule
\end{tabular}
\caption{The tool finds a P-recurrence when one exists, finds an algebraic
equation when one exists, and refuses both when the box is too small.}
\label{tab:controls}
\end{table}

A fifth arm --- the \emph{null control} --- is what makes any positive verdict in
this paper trustworthy: $199$ terms of the by-area king series, the same length
and comparable term size to the perimeter series, come back EXCLUDED with
nullity $0$ at every box where the perimeter series is a candidate. A
$199$-term matrix of that shape does not manufacture rank defects.

\subsection{$\mu$, $\theta$, and the amplitude}

The discriminator for the shape of the correction is the ratio of successive
differences of $r_n = a(n{+}1)/a(n)$: a constant below $1$ means the correction
is geometric --- an isolated subdominant singularity, no power-law factor ---
while drift toward $1$ would mean a branch point and $\theta \ne 0$.

\begin{table}[H]
\centering\small
\begin{tabular}{l l l}
\toprule
series & $d_n/d_{n-1}$ & reading \\
\midrule
king (HV-convex) & $0.48100879371$, constant to all 12 printed digits & geometric \\
control (\oeis{A067675}) & $0.625135968591$, constant to all 12 printed digits & geometric \\
\bottomrule
\end{tabular}
\caption{Measured over the last five $n$ of each $700$-term series.}
\label{tab:theta}
\end{table}

The direct probe agrees: $n(r_n/\mu - 1)$, which tends to $\theta$ under a power
law, sits at the $10^{-216}$ level for the king series and $10^{-137}$ for the
control --- that is, at zero. So $\theta = 0$ for both, there is no power-law
correction, and nothing resembling a stretched exponential:
\[
  a(n) = C\mu^n\bigl(1 + O(\rho^n)\bigr), \qquad
  \rho = 0.481008794 \text{ (king)}, \quad 0.625135969 \text{ (control)}.
\]

Trusted digits are \emph{measured}, not asserted: the whole pipeline is re-run
on the series truncated by $30$ terms, and the reported figure is the minimum
over three independent comparisons, minus a two-digit guard.

\begin{quote}\small
$\mu$, convex king animals by area, $199$ trusted digits:\\[2pt]
\texttt{3.1289432697308862522774479953877541605320912219043941349649746499492}\\
\texttt{4438583718257615823063288164783234634835221018937039608167576493048146}\\
\texttt{233778414064847543298805623983850111109827008627878918922549}
\end{quote}

\begin{quote}\small
$\mu$, convex polyominoes by area (\oeis{A067675}), $121$ trusted digits:\\[2pt]
\texttt{2.3091385933304947310987203050172125319118144725816284016944029002844}\\
\texttt{56440748316842717281615774412174374610237798122137142}
\end{quote}

The control's value is an external check on the whole apparatus: Bender's
published constant is $2.30914$~\cite{bender1974}, and the non-king transfer
matrix plus this extrapolation reproduce it and extend it by $116$ digits.

\subsection{Is $\mu$ algebraic? Not in the searched box}

An integer-relation search on $(1,\mu,\dots,\mu^d)$, run at the trusted
precision only. The discipline that matters: a hit is reported as
REJECTED-ARTIFACT unless $(d{+}1)\log_{10}(\text{height})$ is well below half
the trusted digit count, because PSLQ always returns \emph{something} once the
unknowns can absorb the input precision.

\begin{table}[H]
\centering\small
\begin{tabular}{l r r l l}
\toprule
$\mu$ & digits fed & degree $\le$ & height $\le$ & result \\
\midrule
king & 199 & 20 & $10^{8}$ & none at any degree \\
king & 199 & 30 & $10^{6}$ & none at any degree \\
king & 199 & 12 & $10^{15}$ & none at any degree \\
king & 137 & 20 & $10^{5}$ & none at any degree \\
control & 121 & 12 & $10^{8}$ & none at any degree \\
control & 121 & 20 & $10^{4}$ & none at any degree \\
control & 121 & 20 & $10^{8}$ & hits at degree 15--20, \textbf{all REJECTED-ARTIFACT} \\
\bottomrule
\end{tabular}
\caption{The last row shows what happens when the box outgrows the precision:
capacity $125$--$142$ against $121$ trusted digits. That is the reason the
capacity test exists.}
\label{tab:pslq}
\end{table}

So the falsifiable statement is: \emph{no integer polynomial has $\mu_{\text{king}}$
as a root with degree $\le 20$ and coefficients below $10^8$, nor degree $\le 30$
below $10^6$, nor degree $\le 12$ below $10^{15}$}. A closed form, if one exists,
lies outside those boxes. We do not claim $\mu$ is irrational, and no argument
here bears on it.

\subsubsection*{A worked rejection}

The cubic $3219x^3 - 7630x^2 - 18493x + 33955$ was offered mid-investigation,
produced by feeding a then-published $16$-digit $\mu$ to a computer algebra
service. Its nearest root is
\[
  \begin{aligned}
    &3.12894326973088625330312785623831607097 \\
    \mu = {}&3.12894326973088625227744799538775416053
  \end{aligned}
\]
--- agreement to $18$ digits, disagreement at the $19$th. It is also inside the
PSLQ box already searched (degree $3$, height $3.4\times10^4$), which is why that
search returned nothing at degree $3$.

\section{Perimeter is tame}
\label{sec:perimeter}

By \emph{semiperimeter}
--- for a convex animal the box's $W + H$, with edge-perimeter $2(W{+}H)$ ---
the same class is not merely D-finite but algebraic.

\begin{table}[H]
\centering\small
\begin{tabular}{l l l l}
\toprule
class & statistic & verdict & minimal box \\
\midrule
HV-convex king & area & non-D-finite, non-algebraic & --- \\
HV-convex king & semiperimeter & \textbf{algebraic, degree 2} & $(2,9)$ \\
(dir4, HV-convex) king & area & non-D-finite, non-algebraic & --- \\
(dir4, HV-convex) king & semiperimeter & \textbf{algebraic, degree 4} & $(4,22)$ \\
\bottomrule
\end{tabular}
\caption{Area wild, perimeter tame --- and the wild/tame split survives a
directedness constraint strong enough to cut the population to about $1\%$ by
$s = 200$.}
\label{tab:perim}
\end{table}

The unrestricted quadratic, fitted on $34$ rows and holding on all $166$ later
ones with zero failures over $\Z$:
\[
\begin{aligned}
q_0 &= 2t^2 - 21t^3 + 88t^4 - 196t^5 + 242t^6 - 119t^7 + 72t^8 + 16t^9, \\
q_1 &= -4t + 52t^2 - 252t^3 + 554t^4 - 520t^5 + 96t^6 + 128t^7, \\
q_2 &= 2 - 31t + 176t^2 - 416t^3 + 256t^4 + 256t^5 .
\end{aligned}
\]
$(2,8)$ is EXCLUDED and rational is excluded to $t$-degree $98$, so the box is
minimal in both directions.

The directed quartic has $\deg q_j \le 22$; it was fitted on the first $119$
rows, annihilates all $81$ later ones, and re-substituted into the integer
series annihilates all $200$ rows over $\Z$ with zero failures. $(4,21)$ is
EXCLUDED, as is every algebraic degree below $4$ as far as $199$ terms reach
(degree $3$ to $t$-degree $48$, degree $2$ to $65$, rational to $98$).

\begin{remark}[the coefficients are themselves evidence]
The quartic's coefficients have five digits. A rank defect absorbed from $199$
large terms would produce sixty-digit integers --- and does: the $(19,3)$
P-recurrence for the same object, which is a \emph{non-minimal} description, has
coefficients of about sixty digits.
\end{remark}

\subsection{The nullity law}

The internal consistency check for a positive algebraicity verdict. A genuine
minimal $(K_0,L_0)$ relation forces nullity $(K-K_0+1)(L-L_0+1)$ as the box
opens, since $F^it^jP$ all lie in the larger box. Measured against predicted for
$(K_0,L_0) = (4,22)$, the two grids agree cell for cell over the whole decidable
region; on the null control the same grid is all zeros.

\subsection{Column-convex: solved, and already known}

For contrast, the neighbouring class is exactly solvable.
Temperley's method~\cite{temperley1956} with the last column's height as
catalytic variable --- king contact gives $h + h' + 1$ placements --- closes to a
$2\times2$ linear system and yields
\[
  F(x) \;=\; \frac{x(1-x)^3}{1 - 7x + 13x^2 - 10x^3 + 2x^4},
\]
with growth $4.64468\dots$, a quartic root. Brute-validated to $n \le 8$.

This is a rediscovery: the generating function is that of OEIS \oeis{A187077}
verbatim. Two things about that entry are worth recording, since it carries no
derivation. First, its comment ``Equivalent to a sequence of row-convex
polyhexes (\oeis{A059716})'' is not correct in its plain reading: row-convex
polyhexes are $1,3,11,42,162,\dots$ while this sequence is $1,4,18,83,385,\dots$,
so the two agree only at the first term. Checked by direct enumeration on both
lattices with all fixed polyhexes reproducing \oeis{A001207} as a control. One
placement is the whole difference --- a row of length $h$ admits $h + h' + 1$
positions for a length-$h'$ row beside it on the king lattice against $h + h'$
on the hexagonal one. Second, what the sequence counts is row-convex
\emph{polyplets}; column-convex polyplets are the transpose and give the same
counts.

\section{The subdominant structure}
\label{sec:subdominant}

Proposition~\ref{prop:squeeze} settles the leading term and says nothing about
the corrections, which is where the four-cone class differs from the
unrestricted one. Table~\ref{tab:blocks} locates the corrections too.

Of the four diagonal blocks, exactly two carry exponential weight standalone:

\begin{table}[H]
\centering\small
\begin{tabular}{c l l}
\toprule
block & unrestricted & under the four-cone condition \\
\midrule
$(0,0)$ & sub-exponential (Lemma~\ref{lem:stacks}) & still sub-exponential \\
$(1,0)$ & staircase, growth $\mu = 3.1289\dots$ & untouched \\
$(0,1)$ & mirror staircase, growth $\mu$ & \textbf{truncated, growth $\nu = 2.5146\dots$} \\
$(1,1)$ & sub-exponential (Lemma~\ref{lem:stacks}) & untouched \\
\bottomrule
\end{tabular}
\caption{The truncated $(0,1)$ block is the exception, and the only one.}
\label{tab:subdom}
\end{table}

The truncated $(0,0)$ block stays sub-exponential for a \emph{different} reason
from Lemma~\ref{lem:stacks}'s: with $h' \ge h$, the entry count
$\min(2, h'-h+1)$ is $2$ only at a strict rise, strict rises take distinct
heights, so there are at most $\sqrt{2n}$ of them and the weight is at most
$2^{\sqrt{2n}}$. The truncated $(0,1)$ block is not height-monotone --- $h' \le
h$ at weight $2$, $h' = h+1$ at weight $1$, $h' > h+1$ forbidden --- so heights
can climb one row at a time and Lemma~\ref{lem:stacks} does not apply. Counted
by area it is
\[
  T(n) \;=\; 1,\, 3,\, 8,\, 21,\, 54,\, 138,\, 350,\, 885,\, 2233,\, 5626,\, \dots
\]
with no OEIS match on ten terms. (Its six-term prefix collides with
\oeis{A127358} and two others, all of which diverge by $n = 9$: $2230$ against
$2233$.)

\subsection{The split is exact, and both halves have proved rates}

Partition the four-cone class by whether an animal's phase path ever visits
$(0,1)$:
\[
  D(n) \;=\; D_{\mathrm{asc}}(n) + D_{\mathrm{desc}}(n).
\]

\begin{lemma}
\label{lem:T}
$\lim_n T(n)^{1/n}$ exists. Write $\nu$ for it.
\end{lemma}

\begin{proof}
Let $T_h(n)$ count runs whose first column has height $h$, and write
$T_1 = T_{h=1}$.
(i) $T(n) \ge 2T(n{-}1)$: append a height-$1$ column, which the truncated kernel
admits at weight $2$ from any height.
(ii) $T_h(n) \le 2T(n{-}h)$: delete the first column, the step out of it having
weight at most $2$. With (i), $T_h(n) \le 2^{2-h}T(n{-}1)$, so
$\sum_{h\ge4}T_h(n) \le T(n{-}1)/2 \le T(n)/4$ and
$T(n) \le 4\max_{h\le3}T_h(n)$.
(iii) $T_h(n) \le T_1\bigl(n + h(h{-}1)/2\bigr)$: prepend the climb
$1,2,\dots,h{-}1$, every step of which is the weight-$1$ transition $h' = h+1$;
for $h \le 3$ that costs at most $3$. So $T(n) \le 4T_1(n{+}3)$.
(iv) $T_1(i)T_1(j) \le T_1(i{+}j)$: join two runs that both start at height $1$
by the step ``last height $\to 1$'', weight $2$; the pair is recovered from the
join and the prefix area.
Fekete gives $\lim T_1(n)^{1/n} = \sup_n T_1(n)^{1/n}$, and
$T_1 \le T \le 4T_1(\cdot{+}3)$ squeezes $T$ onto the same limit.
\end{proof}

\begin{proposition}
\label{prop:split}
$\lim_n D_{\mathrm{asc}}(n)^{1/n} = \mu$ and
$\lim_n D_{\mathrm{desc}}(n)^{1/n} = \nu$.
\end{proposition}

\begin{proof}
$D_{\mathrm{asc}} \ge M$ termwise: a staircase animal has $d \ge 0$ and
$d \ge s$ at every step, so it can never enter $(0,1)$, which needs $d \le 0$
and $d < s$. With $D_{\mathrm{asc}} \le D$ and
Proposition~\ref{prop:squeeze}, the first claim follows.

For the second, Lemma~\ref{lem:blocks}'s factorisation with $T$ in place of $M$
gives $D_{\mathrm{desc}}(n) \le 2(n{+}1)^4E(n)^2C_\varepsilon(\nu+\varepsilon)^n$,
so $\limsup D_{\mathrm{desc}}(n)^{1/n} \le \nu$. Conversely take a $T$-run of
area $m$ whose first column is $[0,h{-}1]$ and prepend the column $[0,h]$: the
joining step has $d = 0 \le 0$ and $s = 1 > d$, so it enters $(0,1)$, the run's
own steps keep it there, and the result is a (dir4, HV-convex) animal of area
$m + h + 1$. Restricting to $h = 1$ gives $D_{\mathrm{desc}}(n) \ge T_1(n{-}2)$,
so $\liminf D_{\mathrm{desc}}(n)^{1/n} \ge \nu$ by Lemma~\ref{lem:T}.
\end{proof}

Measured,
\begin{quote}\small
$\nu = 2.5145796438787291885104371943430998201410308559007832719354957107238\\
4099229632690389299337761749953732175539964789209170137762869374741150849461805388232$
\end{quote}
to $153$ trusted digits by the conservative extrapolation and $242$ by Prony's
method on the same terms.

\subsection{The Prony spectrum}

Reading the exponential spectrum of each $700$-term series by Prony's method ---
solve for the constant-coefficient recurrence the tail obeys best and take the
roots of its characteristic polynomial, so every $\lambda_i$ is measured in one
solve, with no peeling and no assumed model:

\begin{table}[H]
\centering\small
\begin{tabular}{l l l l l}
\toprule
series & $\lambda_1$ & $\lambda_2$ & $\lambda_3$ & $\lambda_4$ \\
\midrule
staircase & $3.12894326973088\dots$ & $1.50504922775900\dots$ & $1.28433727098118\dots$ & $-1.25776216033063\dots$ \\
HV-convex & same & same & same & same \\
$D_{\mathrm{asc}}$ & same & same & same & same \\
block $T$ & $2.51457964387872\dots$ & $1.43040460381247\dots$ & $1.24846593371011\dots$ & $1.17441805313317\dots$ \\
$D_{\mathrm{desc}}$ & same & same & same & same \\
$D$ = dir4 & $3.12894326973088\dots$ & $2.51457964387872\dots$ & $1.50504922775900\dots$ & $1.43040460381247\dots$ \\
\bottomrule
\end{tabular}
\caption{$D$'s spectrum is the other two interleaved, and the interleaving keeps
going. Agreements run to $310$, $246$, $236$ digits on the leading claims, always
against a ``bearable'' figure --- the weaker of the two trusted counts. Negative
controls, required to fail and doing so: the block's $\lambda_1$ and $\lambda_2$
match no root of $D_{\mathrm{asc}}$ (best $0$ and $1$ digits), and the
staircase's $\lambda_2$ matches no root of $D_{\mathrm{desc}}$ (best $1$).}
\label{tab:prony}
\end{table}

\subsection{What the split does not buy}

\begin{conjecture}
\label{conj:sharp}
$D_{\mathrm{asc}}(n) = C\mu^n(1 + O(\rho^n))$ with $\rho = 0.48100879370959\dots$
and $D_{\mathrm{desc}}(n) = C'\nu^n(1 + O(\sigma^n))$ with
$\sigma = 0.568844421888\dots$, where $C = 0.4503031857123411823536\dots$ and
$C' = 2.5795006951239199769\dots$
\end{conjecture}

Granting it, $D(n) = C\mu^n + C'\nu^n + O((\mu\rho)^n + (\nu\sigma)^n)$ with
$\mu\rho = 1.50504\dots$ and $\nu\sigma = 1.43040\dots$, both below $\nu$, so the
subdominant exponential of $D$ is exactly $\nu$.

\emph{This is not a gap peculiar to the subdominant.} The repository has no
proof that \emph{any} series in this family has a sharp asymptotic at all:
Proposition~\ref{prop:squeeze} and Lemma~\ref{lem:supermul} give exponential
rates, never $A(n) \sim C\mu^n$, and the amplitudes in
Table~\ref{tab:fourrungs} are all measured, none derived. Asking for the second
exponential as a theorem is asking for strictly more than is known about the
first.

A route exists and it is specific. The \oeis{A225114} entry carries a
\emph{conjectured} continued fraction (Kurkov, September 2024),
\[
  \text{g.f.} = \cfrac{1}{2 - \cfrac{1}{1 - \cfrac{x}{1 - \cfrac{x}{1 - \cfrac{x^2}{1 - \cfrac{x^2}{1 - \cdots}}}}}},
\]
whose depth-$24$ convergent reproduces $41$ terms and whose poles, as $1/x$,
land on the measured staircase spectrum --- $30$, $27$, $14$ and $16$ digits on
$\lambda_1,\dots,\lambda_4$, improving with depth. A proof of that continued
fraction would supply the meromorphic continuation and the spectral gap that a
coefficient-level argument cannot reach.

\subsection{A correction to an earlier reading}

An earlier version of this work read the two series as having \emph{different}
subdominant behaviour, with the four-cone one lacking the $0.4810$. That reading is wrong:

\begin{table}[H]
\centering\small
\begin{tabular}{l l l}
\toprule
series & $d_n/d_{n-1}$ & what it is \\
\midrule
HV-convex, staircase & $0.481008794$ & $\lambda_2/\lambda_1$, the staircase block's own correction \\
$D_{\mathrm{asc}}$ & $\mathbf{0.48100879371}$ & the same one \\
$D_{\mathrm{desc}}$, block $T$ & $0.568844421888$ & the block's own correction \\
(dir4, HV-convex) & $0.803651401483$ & $\nu/\mu$, the \emph{extra} exponential in front \\
\bottomrule
\end{tabular}
\caption{The restricted series has an extra exponential in front of the shared
one, and the $d_n/d_{n-1}$ diagnostic reports only the largest correction
present.}
\label{tab:tableB}
\end{table}

The second row is measured on four-cone-directed animals directly, with no
peeling and no annihilator. A by-product of the split: deleting the
descending half stops the $2.5146$ from polluting the extrapolation, and the
same run's $\mu$ jumps from $51$ trusted digits to $200$.

It remains true that this rules out a bijection between the two classes, and now
for a proved reason rather than an inferred one: a bijection would have to
account for a half of one side, growing at $\nu$, that has no counterpart on the
other.

\section{The amplitude ratio}
\label{sec:amplitude}

The ratio $r = C_{\text{dir4}}/C_{\HV}$ was for a long time just a measured
decimal. It has a closed expression.

\begin{proposition}[the mirror halves]
\label{prop:mirror}
Partition the HV-convex animals by which middle phase their path visits. Then
$A_{(1,0)}(n) = A_{(0,1)}(n)$ for every $n$, and the remainder --- the paths
$(0,0)\to(1,1)$ --- is sub-exponential. Hence if $C_{\HV} = \lim A(n)/\mu^n$
exists, $C_{\HV} = 2\,C(A_{(1,0)})$.
\end{proposition}

\begin{proof}
The vertical mirror $[b(j),t(j)] \mapsto [-t(j),-b(j)]$ is an area-preserving
involution of the class that leaves every column height alone. It carries
$d(j) = b(j{+}1)-b(j)$ to $-(t(j{+}1)-t(j))$, so ``$b$ has risen'' becomes
``$t$ has fallen'': the phase bits swap. Phases $(0,0)$ and $(1,1)$ are fixed
and $(1,0)\leftrightarrow(0,1)$, so the involution is a bijection between the
first two parts. The remainder factors as a phase-$(0,0)$ block joined to a
phase-$(1,1)$ block at one step, both stacks by Lemma~\ref{lem:stacks}, so it is
at most $(n{+}1)^2E(n)^2 = \mu^{o(n)}$.
\end{proof}

\begin{remark}[attribution --- 2 of 3]
\label{rem:attr2}
The mirror equality of the two middle blocks is Gouyou-Beauchamps and Leroux's
$H_{02} = Pa = H_{20}$ in the convex-polyomino setting~\cite{gblRoux2004}; see
Remark~\ref{rem:attr1}. What this proposition adds here is the factor
$\tfrac12$ and the sub-exponential bound on the remainder.
\end{remark}

\begin{proposition}[the staircase eigenvector is a three-term recurrence]
\label{prop:phi}
Let $T(x)_{h,h'} = x^{h'}(\min(h,h')+1)$ for $h,h' \ge 1$ be the staircase
block's operator, and let $\varphi$ solve
\[
  \varphi(0) = 1, \quad \varphi(1) = 2, \quad
  \varphi(h) = \bigl(2 - x^{h-1}\bigr)\varphi(h{-}1) - \varphi(h{-}2)
  \quad (h \ge 2).
\]
For $0 < x < 1$: $T(x)\varphi = \varphi$ if and only if
$\varphi(h) - \varphi(h{-}1) \to 0$; that happens at exactly one $x = x_c$; and
$x_c = 1/\mu$.
\end{proposition}

\begin{proof}
Write $Q(h) = \sum_{h'>h}x^{h'}\varphi(h')$. Splitting $T(x)\varphi = \varphi$
at $h' = h$ gives $\varphi(h) = \sum_{h'\le h}x^{h'}(h'{+}1)\varphi(h') +
(h{+}1)Q(h)$; differencing once gives
$\varphi(h)-\varphi(h{-}1) = x^h\varphi(h) + Q(h)$, and differencing again gives
the recurrence. Conversely, \emph{define}
$Q(h) := \varphi(h)-\varphi(h{-}1)-x^h\varphi(h)$ from the recurrence's
solution; the second difference makes $Q(h{-}1)-Q(h) = x^h\varphi(h)$ an
identity, so $Q(h) = Q(\infty) + \sum_{h'>h}x^{h'}\varphi(h')$ and the
eigenvector equation holds exactly when $Q(\infty) = 0$. Since $x < 1$ the
recurrence's coefficient tends to $2$ and its solution space is spanned by one
asymptotically constant and one asymptotically linear solution; $Q(\infty) = 0$
says the linear one is absent, one analytic condition on $x$.

At such an $x$, $\varphi$ is positive: the differences
$\delta(h) = \varphi(h)-\varphi(h{-}1)$ obey
$\delta(h) = \delta(h{-}1) - x^{h-1}\varphi(h{-}1)$ with $\delta(1) = 1$, so
$\delta$ strictly decreases while $\varphi > 0$, and $\delta(h) \to 0$ forces
$\delta > 0$ throughout --- $\varphi$ increases from $2$ to $2.5374225302\dots$
$T(x)$ is positive and compact on $\{f : |f(h)| \le C(1+h)\}$, so by
Krein--Rutman~\cite{kreinRutman1948} a positive eigenvector's eigenvalue is the
spectral radius. The staircase generating function is
$M(x) = v_0(x)(I-T(x))^{-1}\mathbf 1$ with $v_0(h) = x^h$, of radius $1/\mu$,
and the spectral radius of $T(x)$ increases continuously in $x$, so it reaches
$1$ at $x = 1/\mu$.
\end{proof}

Shooting on that recurrence computes $\mu$ to arbitrary precision in
$O(h_{\max})$ operations, with no series and no extrapolation: it reproduces
all $199$ banked digits, continues past them, and reaches $987$ digits in
$1.6$ seconds on one core, against a $700$-term enumeration plus extrapolation
for $200$.

\begin{proposition}[the ratio]
\label{prop:ratio}
Write $U(h,h') = h'-h+1$ and $U_4(h,h') = \min(2, h'-h+1)$ for $h' \ge h$ (the
$(0,0)$ block, unrestricted and four-cone-truncated), $L(h,h') = \min(h,h'{+}1)$
(the map $(0,0)\to(1,0)$, and also $(0,0)\to(0,1)$ --- the same matrix, which is
Proposition~\ref{prop:mirror} at the operator level), each carrying $x^{h'}$,
and
\[
  w = v_0(I-U)^{-1}L, \qquad w_4 = v_0(I-U_4)^{-1}L,
  \qquad\text{evaluated at } x_c .
\]
If $C(D_{\mathrm{asc}}) = \lim D_{\mathrm{asc}}(n)/\mu^n$ and
$C_{\HV} = \lim A(n)/\mu^n$ exist, then
\begin{equation}
  r \;=\; \frac{C_{\text{dir4}}}{C_{\HV}}
    \;=\; \frac12\cdot\frac{w_4\cdot\varphi}{w\cdot\varphi} .
  \label{eq:ratio}
\end{equation}
\end{proposition}

\begin{proof}
The phase automaton makes the path decomposition exact as generating functions:
$A_{(1,0)}(x) = w(x)(I-T(x))^{-1}q(x)$ with
$q = 1 + L_{(1,0)\to(1,1)}(I-T_{(1,1)})^{-1}\mathbf 1$, and
$D_{(1,0)}(x) = w_4(x)(I-T(x))^{-1}q(x)$ with the \emph{same} $T$ and the
\emph{same} $q$ --- the four-cone floor bites only inside the $(0,0)$ block,
since every step out of $(1,0)$ and every step within $(1,1)$ already has
$d \ge 0$. The $(0,0)$ and $(1,1)$ blocks are stacks
(Lemma~\ref{lem:stacks}), so $w$, $w_4$ and $q$ are finite at $x_c$. By
Proposition~\ref{prop:phi} and Krein--Rutman the pole of $(I-T(x))^{-1}$ at
$x_c$ is simple with rank-one residue $\varphi\psi^{\mathsf T}/\langle\psi,
\varphi\rangle$, so the two limits
$\lim_{x\to x_c^-}(1-\mu x)A_{(1,0)}(x)$ and
$\lim_{x\to x_c^-}(1-\mu x)D_{(1,0)}(x)$ are $N(w\cdot\varphi)$ and
$N(w_4\cdot\varphi)$ with the same $N$, which cancels. If the amplitudes exist
then Abel's theorem identifies those limits with $C(A_{(1,0)})$ and
$C(D_{(1,0)})$. Finally $C_{\HV} = 2C(A_{(1,0)})$
(Proposition~\ref{prop:mirror}) and $C_{\text{dir4}} = C(D_{(1,0)})$, the
$D_{\mathrm{desc}}$ and ``via neither'' parts being $o(\mu^n)$ by
Proposition~\ref{prop:split} and Lemma~\ref{lem:stacks}.
\end{proof}

The hypothesis is strictly weaker than Conjecture~\ref{conj:sharp}, which asks
for the error terms as well. What \eqref{eq:ratio} settles is what a closed form
would have to look like: a ratio of two $q$-series in $q = 1/\mu$, both computed
by $O(h_{\max})$ prefix-sum recurrences, with
$w\cdot\varphi = 1.29770192341040021939895\dots$ and
$w_4\cdot\varphi = 1.19935960397018718067118\dots$

\subsection{The number, and three determinations of it}

\begin{quote}\small
$r$, $251$ trusted digits:\\[2pt]
\texttt{0.4621090492099424400356623877003058328411634397299804247065092921445}\\
\texttt{22950511444483541019804848618234062089755257041466533204384387638170139}\\
\texttt{068144091014356289934131553812178862045404673728300959682539208197483810}\\
\texttt{09549588646431928574874902351771539511812}
\end{quote}

\begin{table}[H]
\centering\small
\begin{tabular}{l l r}
\toprule
route & own cross-check & trusted \\
\midrule
$D/A$ extrapolated & raw vs Aitken, $n=600$ vs $700$ & 54 \\
$D_{\mathrm{asc}}/A$ extrapolated & raw vs Aitken, $n=600$ vs $700$ & 185 \\
identity~\eqref{eq:ratio} & $h_{\max}$ $1200$ vs $1000$ & 492 \\
\bottomrule
\end{tabular}
\caption{The extrapolation and the identity agree to $293$ digits. The two
routes share no arithmetic: the identity comes from an exact operator
factorisation and the extrapolation touches no operator at all. The first row
reproduces the historical $54$-digit figure character for character, so it
doubles as the control on the method --- the $185$ is bought by the phase split
and by nothing else.}
\label{tab:ratio}
\end{table}

Negative controls, all required to fail and doing so: dropping the factor
$\tfrac12$ agrees with the measurement to $0$ digits; truncating the $(0,0)$
block at $\min(3,\cdot)$ instead of $\min(2,\cdot)$ agrees to $1$; and the
$\varphi$ recurrence at $x = 1/3.13$ fails its own eigen-equation at $10^{-2}$
where $x_c$ gives $10^{-528}$.

\subsection{Is $r$ algebraic?}

At $251$ digits, no integer relation in any in-capacity box:
\[
  \text{degree} \le 30 \text{ at height} \le 10^4, \;\le 20 \text{ at } 10^5,
  \;\le 10 \text{ at } 10^{11}, \;\le 5 \text{ at } 10^{20},
  \;\le 3 \text{ at } 10^{30}, \;\le 2 \text{ at } 10^{40}.
\]
Nor over $\Q(\mu)$, the natural field since $\mu$ is the constant both series
share: $\mu$-degree $\le 5$ with $r$-degree $\le 6$ at height $\le 10^2$, and
three further boxes. A positive control runs on every pass --- $(1+\sqrt2)/3$
must be found at degree $2$, and is.

$\mu$'s own arithmetic status is what limits how much this can mean. There is no
proof that $\mu$ is irrational, let alone transcendental, so $\Q(\mu)$ is a
field of unknown degree and ``algebraic over $\Q(\mu)$'' cannot be settled
either way by any argument available here. What \eqref{eq:ratio} contributes is
a reason to expect the answer to be no: $r$ is a ratio of two $q$-series
evaluated at $q = 1/\mu$, and $\mu$ is itself the reciprocal of the zero of an
entire $q$-function (Proposition~\ref{prop:phi}) --- the same shape as Kurkov's
conjectured continued fraction.

\section{Relation to existing work}
\label{sec:related}

\paragraph{Attribution --- 3 of 3.} The block decomposition of
Lemma~\ref{lem:blocks}, the identification of the extreme blocks as stacks and
the middle blocks as staircases, and the mirror equality of
Proposition~\ref{prop:mirror} are all present in Gouyou-Beauchamps and
Leroux~\cite{gblRoux2004}, \S2.3, for convex polyominoes on the honeycomb
lattice. Their square-lattice companion~\cite{lrr1998} reaches the same objects
by a different route --- partitions, stacks, shifted stacks and directed convex,
via Temperley and Bousquet-M\'elou --- so~\cite{gblRoux2004} is the source to
cite for the decomposition itself.

What is not theirs, and is what this paper claims: the king lattice, where the
middle kernel is $\min(h,h')+1$ rather than $\min(h,h')$ and neither class is
exactly solved; the squeeze over \emph{every} intermediate class
(Proposition~\ref{prop:squeeze}); the exponential split and its two proved rates
(Lemma~\ref{lem:T}, Proposition~\ref{prop:split}); the eigenvector recurrence
and the shooting method (Proposition~\ref{prop:phi}); and the amplitude-ratio
identity (Proposition~\ref{prop:ratio}).

\paragraph{The square-lattice analogue.} Bender's $2.30914$~\cite{bender1974} for
convex polyominoes by area, and the parallelogram subclass \oeis{A006958}
sharing it (measured here at $2.309138593330495$, flat from $n = 100$ to
$n = 400$), mean that Proposition~\ref{prop:squeeze}'s \emph{conclusion} could
have been read off the solved models one lattice over. It is the generality ---
every intermediate class, on a lattice where nothing is solved --- that is the
contribution.

\paragraph{Perimeter.} Delest and Viennot~\cite{delestViennot1984} is the
classical algebraicity result and is used here as the positive control for the
algebraicity guesser. \S\ref{sec:perimeter}'s finding is the king analogue with
an explicit minimal box.

\paragraph{Column-convex.} \S\ref{sec:perimeter} records a rediscovery of
\oeis{A187077} by Temperley's method~\cite{temperley1956}, together with a
correction to that entry's comment.

\section{Open problems}

\begin{openproblem}[a sharp asymptotic, for anything in this family]
Prove $A(n) \sim C\mu^n$, or the same for $M(n)$. Nothing in this paper reaches
it, and Conjecture~\ref{conj:sharp} depends on it. The most concrete route is a
proof of Kurkov's conjectured continued fraction for \oeis{A225114}, which would
supply the meromorphic continuation and the spectral gap.
\end{openproblem}

\begin{openproblem}[non-D-finiteness as a theorem]
Theorem~\ref{thm:excl} is an exclusion on $700$ terms, not a proof. The
arithmetic obstruction that works for the height-anisotropic king generating
function does not obviously apply here, since the natural slicing of a convex
class is not by height.
\end{openproblem}

\begin{openproblem}[the block series $T$]
$1, 3, 8, 21, 54, 138, 350, 885, 2233, 5626, \dots$ has no OEIS match and no
conjectured form. It is the truncated $(0,1)$ block, its growth constant $\nu$
is the subdominant exponential of the four-cone class, and it deserves the same
treatment \oeis{A225114} has had.
\end{openproblem}

\section{Reproduction}
\label{sec:repro}

\begin{table}[H]
\centering\small
\begin{tabular}{l l}
\toprule
claim & check \\
\midrule
$A(n)$ to $n = 700$ & \texttt{make build/convex\_area\_tm \&\& build/convex\_area\_tm 700 1} \\
the control to $n = 700$ & \texttt{build/convex\_area\_tm 700 0} \\
Theorem~\ref{thm:excl} & \texttt{make build/prec\_guess}; \texttt{build/prec\_guess prec ... 24 24} \\
Table~\ref{tab:controls} & \texttt{make gate-convex-dfinite} \\
$\mu$, $\theta$, the amplitudes & \repo{experiments/convex\_growth.py} \\
Table~\ref{tab:pslq} & \repo{experiments/convex\_growth.py} \texttt{--algdeg} \\
Lemma~\ref{lem:supermul}, measured & \repo{experiments/staircase\_supermul.py} \\
Lemma~\ref{lem:supermul}, in Lean & \texttt{cd polyplets \&\& lake build} \\
Lemma~\ref{lem:blocks}, brute force & \repo{experiments/monotone\_block\_growth.py} \\
the block series $T$ & \repo{experiments/dir4\_descent\_block.py}, \repo{experiments/descent\_block\_oracle.py} \\
Table~\ref{tab:prony} & \repo{experiments/prony\_spectrum.py} \\
Propositions~\ref{prop:phi}--\ref{prop:ratio} & \repo{experiments/amplitude\_feed\_vectors.py} \\
the PSLQ boxes for $r$ & \repo{experiments/amplitude\_pslq.py} \\
\S\ref{sec:perimeter} & \repo{experiments/dir4\_perim\_find\_alg.py}, \texttt{make gate-dir4-perim-alg} \\
column-convex & \repo{experiments/colconvex\_king.py} \\
the squeeze, pinned & \texttt{make gate-middle-kingdom} \\
\bottomrule
\end{tabular}
\caption{Where each claim is checked. Every proof-relevant item is either a
paper proof above or a Lean theorem; everything in this table is a measurement
or a reproduction of one.}
\label{tab:repro}
\end{table}

\bibliographystyle{plain}
\bibliography{shared/refs}

\end{document}
